% Source PDF: source/普通高考/2010/2010安徽文.pdf
% Page: 2; no embedded bitmap matched because the source contains the data prompt only.
% Grid evidence: tmp/repaint_2010_anhui_liberal/grids/page-02-grid100.jpg.
% Elements: seven contiguous class bars for pollution-index intervals [40,50),...,[100,110), axes, ticks, and labels.
% Relations: class width is 10; frequencies from the 30 listed observations are 2,1,3,7,10,5,2,
% so the bar heights (frequency/group width) are 1/150,1/300,1/100,7/300,1/30,1/60,1/150.
% solid segments: x/y axes and seven bar outlines; dashed segments: none.
% labels: the y-axis is frequency/group width and the x-axis is pollution index; ticks stay clear of bars.
\begin{tikzpicture}
\begin{axis}[
    width=10.2cm,
    height=5.8cm,
    axis lines=left,
    axis line style={->,black},
    xmin=39,
    xmax=111,
    ymin=0,
    ymax=0.040,
    xtick={40,50,60,70,80,90,100,110},
    ytick={0,0.01,0.02,0.03,0.04},
    yticklabel=\empty,
    scaled y ticks=false,
    xlabel={污染指数},
    ylabel={},
    xlabel style={at={(axis description cs:1.02,0)},anchor=west,inner sep=0pt},
    tick style={black},
    clip=false,
    enlargelimits=false,
]
    \addplot[
        ybar interval,
        draw=black,
        fill=white,
        line width=0.75pt,
    ] coordinates {
        (40,0.0066667)
        (50,0.0033333)
        (60,0.0100000)
        (70,0.0233333)
        (80,0.0333333)
        (90,0.0166667)
        (100,0.0066667)
        (110,0)
    };
    \node[anchor=south west,inner sep=0pt,font=\normalsize]
        at (axis cs:40,0.0385) {\(\dfrac{\text{频率}}{\text{组距}}\)};
    \node[anchor=north east,inner sep=1pt,font=\normalsize,xshift=-5pt]
        at (axis cs:39,0) {0};
    \frequencyylabel[-4pt]{39}{0.0100000}{0.0100}
    \frequencyylabel[-4pt]{39}{0.0200000}{0.0200}
    \frequencyylabel[-4pt]{39}{0.0300000}{0.0300}
    \frequencyylabel[-4pt]{39}{0.0400000}{0.0400}
\end{axis}
\end{tikzpicture}
